\section{Concrete finite arithmetic sentences}
\label{sec:contracts}

This section fixes representative internal statements rather than leaving
``formalize the analytic theorem'' as an unspecified goal. Integer and
rational coding conventions can be chosen freely, provided their
operations are compiled to the repository's primitive-recursive syntax.
Invalid numerical codes are assigned fixed default values. A conditional
Boolean checker returns one on an invalid hypothesis code; a constructor's
separate validity theorem ensures that its own outputs are genuine
certificates.

\subsection{The primary counting sentence}
Fix the deterministic height selector and contour constructor of
Section~\ref{sec:regular-heights}. For integer $n\ge8$ let
$\mathsf{Gen}(n)$ be its complete rational certificate and
$C(n)$ its crossing-sum integer. Return zero for $n<8$ in the
definition of $C$. Let $\mathsf{Valid}$ recompute the specified xi
approximants, verify the contour mesh, prescribed error bounds and
positive sample margins, and recompute the crossing sum. Two separate
internal theorems are required:
\begin{equation}
 \ipra\vdash\forall n\ge8\quad
            \mathsf{Valid}(\mathsf{Gen}(n))=1,
 \label{eq:generator-sentence}
\end{equation}
and the numerical estimate below. The first is supplied by the
minimum-modulus and finite pigeonhole argument, including the other
three edges of the rectangle.

Let $m(n,p)$ and $\ell(n,p)$ be the fixed rational approximants, with
error at most $2^{-p}$, to
\[
 \mathcal M(n)=\frac{n}{2\pi}\log\frac{n}{2\pi}-\frac{n}{2\pi},
 \qquad \log(n+5),
\]
respectively. Their definitions use the log and exponential routines
and extra working precision sufficient for multiplication by $n$ and
division by $2\pi\ge6$. Define a Boolean function
$\mathsf{CountCheck}(n,p)$ to return one if $n<8$, and otherwise to
test the exact rational inequality
\begin{equation}
 |C(n)-m(n,p)|\le9\ell(n,p)+40+10\,2^{-p}.
 \label{eq:count-matrix}
\end{equation}
Then the intended object sentence is precisely
\begin{equation}
       \ipra\vdash\forall n\,\forall p\quad
                   \mathsf{CountCheck}(n,p)=1.
 \label{eq:count-sentence}
\end{equation}
The $10\,2^{-p}$ budget is the sum of the error in $m$ and nine
times the error in $\ell$. The proof of (RVM-PR) establishes this
inequality with $n,p$ free. For a supplied regular contour, the analogous
checker uses $8\log(U+4)+40$ and a $9\,2^{-p}$ budget.

\subsection{Representative zero-exclusion and identity sentences}
For integer $H\ge1$ and a rational input $s$ in the rectangle
$|\Im s|\le H$, $|\Re s-1|\le\varepsilon_H$, the lower bound
$b_H$ in Theorem~\ref{thm:elem-F-clearance} is rational. Compute the
symbolically pole-removed numerator $F$ with error at most $b_H/8$.
The returned rational $v$ must satisfy
\[
                         |v|^2\ge(7b_H/8)^2.
\]
The checker tests the input inequalities, recomputes $v$, and performs
this squared-norm comparison. The working precision is the bounded
binary logarithm of $8/b_H$ plus the stated arithmetic error budget.
Thus this is one uniform finite test at an explicit precision, rather
than a negation of an undecidable equality.

For $\xi$ symmetry, use one common box and the same theta summation
parameters at $s,1-s,\bar s$. The rational expression for the completion
then satisfies the reflection identity exactly. The conjugation identity
is also exact if the elementary evaluator treats conjugate inputs by
conjugating its rational Taylor polynomial and the chosen real constants.
The Mellin identification remains a separate, substantive finite theorem.

For trivial zeros and their derivative formula, compile the finite gamma
shift identity and the Taylor coefficient identity at the integer input
$-2n$. At precision $p$ the equality checker compares the two prescribed
rational approximants with the sum of their proved error budgets. The
strict simplicity assertion uses the rational lower bound
$(2n)!/(2\,8^{2n})$ to select a fixed successful derivative precision.
No equality test on an arbitrary real or complex input is part of this
checker.

\subsection{Reciprocal sums as Cauchy names}
For each integer $m\ge2$, define $A_m(p)$ by a weighted finite contour
integral at a selected regular height above $T_{p+3}$, using both
upper and lower rectangles. Compute the quadrature and function values
to total Euclidean error at most $2^{-p-3}$. The tail has the same
bound by~\eqref{eq:reciprocal-modulus}. After conversion to the
rational L1 norm, the error is at most $2^{-p-1}$. The internal
statement can therefore be the exact rational inequality
\begin{equation}
 \normone{A_m(p)-A_m(q)}\le2^{-p}+2^{-q}
       \qquad(m\ge2).
 \label{eq:cauchy-sentence}
\end{equation}
All loop bounds remain primitive recursive when $m$ is a numerical
parameter: the weight and its derivatives involve finite powers, and
the lower contour has the fixed rational height $\eta=1/4096$.
For the paired reciprocal sum use its upper-contour weight
$1/[s(1-s)]$ instead. The same sentence and budgets apply.

\subsection{Why these are uniform proofs}
The proof uses several nested loops, but every nesting has a numerical
bound supplied before execution. Examples are the degree-bounded gcd,
finite determinant permutation list, bounded algebraic root grid,
fixed-precision bisections, Euler--Maclaurin truncations, candidate-height
list, contour mesh and crossing sum. Their state is a finite numerical
array. At each loop, all later precision demands are already numerical
parameters. Identities and rational error inequalities about that state
are quantifier-free invariants.

For a bounded family of such inequalities, use the Boolean product of
their characteristic functions. Induction applies to the equation that
this product equals one, with the bound as a free parameter. It does
not apply to a statement quantifying over every precision or asserting
existence of an unspecified real object. Polynomial root names in the
finite algebraic intermediate calculus have one fixed bounded
evaluation routine on their finite codes, so no universal evaluator for
arbitrary primitive recursive programs is being assumed.

These details specify the arithmetic claims and the human proofs of
their uniform bounds. The remaining Lean implementation work is to
compile the fixed programs, prove the indicated invariants in the deeply
embedded theory, and build the actual derivations of
\eqref{eq:generator-sentence}, \eqref{eq:count-sentence}, and
\eqref{eq:cauchy-sentence}. A theorem about their truth in ambient Lean's
standard natural numbers is a different endpoint.
